% vim: textwidth=78 ai spell spelllang=en filetype=tex sw=2

\chapter{Symbiotic}\label{sec:fmics}
In the previous chapter, we discussed the abstract interpretation and its
application to code analysis. We especially aimed at checking code using
finite state automata (FSA).  This formalism is simple and still flexible
enough to describe many often studied program properties including locking
policy in concurrent programs, null-pointer dereferences, resource
allocations, and resource leaks. FSA specification is therefore used in many
static program analysis tools like \textsc{xgcc}~\cite{metal},
\textsc{Slam}~\cite{slam_decade, slam2}, \textsc{Sdv}~\cite{sdv},
\textsc{Blast}~\cite{blast}, \textsc{Esp}~\cite{esp}, or \stan~\cite{stanse}.
All the mentioned tools produce \emph{false positives}, i.e.\ they report
errors that do not correspond to any real error. We showed an example of such
a false positive in the previous chapter. This lead us to use of the symbolic
execution, which we extensively described in Section~\ref{sec:se}. But
the symbolic execution is known to suffer from slowness on large programs or
programs with loops. So we had to overcome this issue. The idea is that we can
find suspicious places in analyzed code by cheap analyses like abstract
interpretation. Then we can take a symbolic executor and verify the problem
exists at those places. And since we know places in the code which we want to
reach, we can use the program slicing to eliminate code that does not affect
the places. We described the program slicing in Section~\ref{sec:slicing}.

We can extend the principle further to instrument the code with some behaviour
and slice with the behaviour as slicing criteria. For example, the behaviour
can be automaton transitions at the same places as would be fired in the
abstract interpretation. So this chapter describes the new approach based on
three well-known methods: instrumentation, slicing and symbolic execution. We
decided to dub the tool implementing this combination \symbio. Exactly because
it is a symbiosis of techniques.  This work is based on our paper published in
Springer's LNCS~\cite{symbiotic} and presented on the 17th International
Workshop on Formal Methods for Industrial Critical Systems 2012.

\paragraph{Chapter Organization}
We structured this chapter as follows. We first describe an example in
Section~\ref{sec:fmics_example} which we use for demonstration of the
technique throughout the whole chapter. Thereafter we describe how we use
the instrumentation (Section~\ref{sec:fmics_instr}), slicing
(Section~\ref{sec:fmics_slicing}), and symbolic execution
(Section~\ref{sec:fmics_symexe}) for error-finding. Then we discuss the
process of implementation and experimental results in
Section~\ref{sec:fmics_impl}. Finally, we suggest some future directions in
Section~\ref{sec:fmics_future} and conclude in
Section~\ref{sec:fmics_conclusion}.

\section{Running Example}\label{sec:fmics_example}
We now roughly sum up the basic principle of a static analysis implemented
in \textsc{xgcc}, \textsc{Esp}, and \textsc{Stanse} and show an example where
it does not work. The example will accompany us through the following sections
to demonstrate each part of the technique. So let us consider the automaton
$\SA(x)$ of Figure~\ref{fig:sm}.
\begin{figure}[htb]
  \centering
    \tikzstyle{loc} = [circle,thick,draw,minimum size=8mm]
    \tikzstyle{err} = [loc, fill=red!10, draw=red]
    \tikzstyle{pre} = [<-, shorten <=1pt, >=stealth', semithick]
    \tikzstyle{post} = [pre, ->]
    \begin{tikzpicture}[node distance=8mm]
        \node (S) {};
        \node [loc] (U) [right=6mm of S] {\code{U}}
          edge [pre] node {} (S);
        \node [loc] (L) [right=18mm of U] {\code{L}}
          edge [pre,bend left=20] node [label=below:\code{lock($x$)}] {} (U)
          edge [post,bend right=20] node [label=above:\code{unlock($x$)}] {} (U);
        \node [err] (DU) [below=of U] {\code{DU}}
          edge [pre] node [label=left:\code{unlock($x$)}] {} (U);
        \node [err] (DL) [below=of L] {\code{DL}}
          edge [pre,near start] node [label=left:\code{lock($x$)}] {} (L);
        \node [err] (RL) [right=6mm of DL] {\code{RL}}
          edge [pre,bend right=10,near start] node [label=above right:\code{return}] {} (L);
    \end{tikzpicture}
    \caption{The state automaton $\SA(x)$ describing errors in manipulation
    with the lock $x$. The nodes \code{U} and \code{L} refer to states
    \emph{unlocked} and \emph{locked}, respectively. The other three nodes
    refer to error states: \code{DU} to \emph{double unlock}, \code{DL} to
    \emph{double lock}, and \code{RL} to \emph{return in locked state}. The
    initial node is \code{U}.}
  \label{fig:sm}
\end{figure}
It describes a lock manipulation including malign transitions. Intuitively,
the automaton represents possible courses of states of a lock referenced by
$x$ along an execution of a program.  The state of the lock is changed
according to the program behaviour.  Whenever the program contains a statement
syntactically subsuming the label of a transition, the transition is fired in
the state automaton. We would like to decide whether there exists any program
execution where an instance of the automaton $\SA(x)$ reaches an error state
for some lock in the program.  Unfortunately, this is not feasible due to
potentially unbounded number of executions and unbounded execution length.
Hence, static analysis tools over-approximate the set of reachable automaton
states.

Let us assume that we want to check the program of Figure~\ref{fig:copy} for
errors specified by the automaton $\SA(x)$.
\begin{figure}[htb]
  \begin{center}
  \begin{tabular}{c}
    \begin{lstlisting}
char *copy(char *dst, char *src, int n, int *L) {(*@\label{line:copy_start}@*)
  int i, len;                            // {U}
  len = 0;                               // {U}(*@\label{line:copy_len0}@*)
  if (src != NULL && dst != NULL) {      // {U}(*@\label{line:copy_if1}@*)
    len = n;                             // {U}
    lock(L);                             // {L}(*@\label{line:copy_lock}@*)
  }                                      // {U,L}
  i = 0;                                 // {U,L}(*@\label{line:copy_i0}@*)
  while (i < len) {                      // {U,L}(*@\label{line:copy_loop}@*)
    dst[i] = src[i];                     // {U,L}
    i++;                                 // {U,L}
  }                                      // {U,L}(*@\label{line:copy_loopE}@*)
  if (len > 0) {                         // {U,L}(*@\label{line:copy_if2}@*)
    unlock(L);                           // {DU,U}(*@\label{line:copy_unlock}@*)
  }                                      // {U,L}(*@\label{line:copy_if2_e}@*)
  return dst;                            // {U,RL}(*@\label{line:copy_ret}@*)
}
    \end{lstlisting}
  \end{tabular}
  \end{center}
    \caption{Function \code{copy} copying a source string \code{src}
      into a buffer \code{dst} using a lock \code{L} to prevent parallel
      writes.}
    \label{fig:copy}
\end{figure}
First, we find all locks in the program and to each lock we assign an instance
of the automaton. In our case, there is only one lock pointed to by \code{L}
and thus only one instance $\SA(\var{L})$. For each program location, we
compute a set over-approximating possible states of $\SA(\var{L})$ after
executions leading to the location. Roughly speaking, we initialize the set in
the initial location to $\{\var{U}\}$ and the other sets to $\emptyset$. Then
we repeatedly update the sets according to the effect of individual program
statements until the fixed point is reached.  The resulting sets for the
program of Figure~\ref{fig:copy} are written directly in the code listing as
comments on the right side.

As we can see, the sets contain two error states: \emph{double unlock} after
the \code{unlock(L)} statement and \emph{return in locked state} in the
terminal location. If we analyze the computation of the sets, we can see
that the first error corresponds to executions going through all lines
\ref{line:copy_start}--\ref{line:copy_if1} and \ref{line:copy_i0}, then
iterating the \code{while}-loop and finally passing lines~\ref{line:copy_if2}
and \ref{line:copy_unlock}.  These execution paths are not feasible due to the
value of \code{len}, which is set to zero on line~\ref{line:copy_len0} and
assumed to satisfy \code{len > 0} on line~\ref{line:copy_if2}. Hence, the
first error is a \emph{false positive}. The second error corresponds to
executions passing all lines \ref{line:copy_start}--\ref{line:copy_i0}, then
iterating the \code{while}-loop and finally going through
lines~\ref{line:copy_if2} and \ref{line:copy_ret}. All these paths are also
infeasible except the one that performs zero iterations of the
\code{while}-loop, which is the only real execution leading to the only real
locking error in the program.

To sum up, static analysis tools like \textsc{xgcc}, \textsc{ESP}, and \stan
are highly flexible, fast and thus applicable to extremely
large software projects like the Linux kernel. The tools examine whole code
and find many error reports. Unfortunately, many of the reports are false
positives. So the tools have to use many techniques for partial elimination of
false positives.  The main source of false positives is related to the fact
that the analysis does not work with data values. In particular, the analysis
does not track connections between variable values and states of automata.
Therefore, when a branching statement is processed, the analysis cannot decide
which states of automata are reachable on the positive branch and which on the
negative one.  It may be illustrated by the \emph{double unlock} false
positive detected in the function of Figure~\ref{fig:copy} because the
analysis does not know that the condition on line~\ref{line:copy_if2} holds
only if the automaton $\SA(\var{L})$ is in state $L$. As manual sorting of
error reports containing a pile of false positives is tedious work, the
practical applicability of such tools is limited.

\section{Our Contribution}

In contrast to the static analysis, test-generation tools based on the
symbolic execution do not suffer from false positives, since the checked
program is actually executed (but on symbolic input instead of concrete one).
A disadvantage of these tools is that they usually detect only low-level
errors representing violations of programming language semantics, i.e.\ 
various types of undefined behavior or crashes. To detect violations of other
program properties like locking policy, the program has to be modified such
that the errors can be detected during the execution. This can be achieved,
for example, by introducing a couple of \code{assert()} statements to specific
program locations.  Another and even more important disadvantage of the tools
is extreme computation cost of the symbolic execution. In particular, programs
containing loops or recursion have typically large or even infinite number of
execution paths and cannot be entirely analysed by the symbolic execution.

\subsection*{New Technique}
In this chapter, we introduce a new fully automatic program analysis technique
offering flexibility of FSA property specification with zero false positive
rate of the symbolic execution. The technique symbolically executes only parts
of the analysed program having impact on the checked property. The basic idea
is very simple:
\begin{enumerate}
\item We get automata describing some program properties. We instrument a
  given program with code tracking behaviour of the automata.
\item The instrumented program is then reduced using a method called
  slicing~\cite{slicing}. The sliced program has to meet the criterion
  to be equivalent to the instrumented program with respect to reachability
  of error states of tracked state machines. Note that slicing may remove
  large portions of the code, including loops and function calls. Hence, an
  original program with an infinite number of execution paths may be reduced
  to a program with a finite number of execution paths.
\item Finally, we execute the sliced program symbolically to find violations
  of the checked property.
\end{enumerate}
Our technique may be used in two ways according to the applied symbolic
execution tool. If we apply a symbolic executor that prefers to explore more
parts of the code (for example by exploring each program loop at most twice),
we may use the technique as a general error-finding technique reporting only
real errors.  On the contrary, if we use a symbolic executor exploring all
execution paths, we may use our technique for classification of error reports
produced by other tools (e.g.\ \textsc{xgcc} or \stan). For each such an error
report, we may instrument the corresponding code only with the state machine
describing the reported error. If our technique finds the same error, it is a
real one.  If our technique explores all execution paths of the sliced code
without detecting the error, it is a false positive. If our technique runs out
of resources, we cannot decide whether the error is a real one or just a false
positive.

We have developed an experimental tool \symbio. It instruments a program with
an automaton describing locking errors (we use a single-purpose
instrumentation so far), then it applies an inter-procedural slicing to the
instrumented code, and it passes the sliced code to symbolic executor
\klee~\cite{klee}. Our experimental results indicate that the
technique can indeed classify error reports produced by \stan applied to the
Linux kernel.

We emphasize the synergy of the three known methods combined in the
presented technique.
\begin{itemize}
\item The instrumentation of a program with code emulating automata provides us
with simple slicing criteria: we want to preserve values of memory places
representing states of automata. Hence, the sliced program contains only the
code relevant to the considered errors specified by automata.
\item The slicing may substantially reduce the code size, which in turn may
remarkably improve performance of the symbolic execution.
\item The application of the symbolic execution brings us another benefit.
While in the standard static analysis, the automata are associated to
syntactic objects (e.g.\ lock variables appearing in a program), we may
associate automata to actual values of these objects. This leads to a higher
precision of error detection.
\end{itemize}

%{{{ Obrazek nekonecneho stromu symbolicke exekuce na neproslicovanem kodu

% \begin{figure}[!htb]
%     \centering
%     \tikzstyle{loc} = [circle,thick,draw,inner sep=0.9pt]
%     \tikzstyle{loc2} = [circle,thick,draw,inner sep=0.1pt]
%     \tikzstyle{pre} = [<-,shorten <=1pt,>=stealth',semithick]
%     \tikzstyle{post} = [->,shorten <=1pt,>=stealth',semithick]
%     \footnotesize
%     \begin{tikzpicture}[node distance=0.65cm]
%         \node [loc] (1) {$2$};
%         \node [loc] (2) [below of=1] {$3$}
%             edge [pre] node {} (1);

%         % ->false branch
%         \node [loc] (3) [left of=2,below of=2] {$7$}
%             edge [pre] node [label=above:F~~] {} (2);
%         \node [loc] (4) [below of=3] {$8$}
%             edge [pre] node {} (3);
%         \node [loc2] (5) [below of=4] {$12$}
%             edge [pre] node {} (4);
%         \node [loc2] (6) [below of=5] {$14$}
%             edge [pre] node {} (5);

%         % ->true branch
%         \node [loc] (7) [right of=2,below of=2] {$4$}
%             edge [pre] node [label=above:~~T] {} (2);
%         \node [loc] (8) [below of=7] {$5$}
%             edge [pre] node {} (7);
%         \node [loc] (9) [below of=8] {$7$}
%             edge [pre] node {} (8);
%         \node [loc] (10) [below of=9] {$8$}
%             edge [pre] node {} (9);

%         % ->true->false branch
%         \node [loc2] (11) [left of=10,below of=10] {$12$}
%             edge [pre] node [label=above:F~~] {} (10);
%         \node [loc2] (12) [below of=11] {$14$}
%             edge [pre] node {} (11);

%         % ->true->true branch
%         \node [loc] (13) [right of=10,below of=10] {$9$}
%             edge [pre] node [label=above:~~T] {} (10);
%         \node [loc2] (14) [below of=13] {$10$}
%             edge [pre] node {} (13);
%         \node [loc] (15) [below of=14] {$8$}
%             edge [pre] node {} (14);

%         % ->true->true->false branch
%         \node [loc2] (16) [left of=15,below of=15] {$12$}
%             edge [pre] node [label=above:F~~] {} (15);
%         \node [loc2] (17) [below of=16] {$13$}
%             edge [pre] node {} (16);
%         \node [loc2] (18) [below of=17] {$14$}
%             edge [pre] node {} (17);

%         % ->true->true->true branch
%         \node [loc] (19) [right of=15,below of=15] {$9$}
%             edge [pre] node [label=above:~~T] {} (15);
%         \node [loc2] (20) [below of=19] {$10$}
%             edge [pre] node {} (19);
%         \node [loc] (21) [below of=20] {$8$}
%             edge [pre] node {} (20);

%         % ->true->true->false->false branch
%         \node [loc2] (22) [left of=21,below of=21] {$12$}
%             edge [pre] node [label=above:F~~] {} (21);
%         \node [loc2] (23) [below of=22] {$13$}
%             edge [pre] node {} (22);
%         \node [loc2] (24) [below of=23] {$14$}
%             edge [pre] node {} (23);

%         % ->true->true->true->true branch
%         \node [loc] (25) [right of=21,below of=21] {$9$}
%             edge [pre] node [label=above:~~T] {} (21);
%         \node [loc2] (26) [below of=25] {$10$}
%             edge [pre] node {} (25);
%         \node [loc] (27) [below of=26] {$8$}
%             edge [pre] node {} (26);
%         \node [] (FINF) [below of=27,left of=27] {}
%             edge [dotted,thick] node {} (27);
%         \node [] (TINF) [below of=27,right of=27] {}
%             edge [dotted,thick] node {} (27);
%     \end{tikzpicture}
%     \caption{Symbolic execution tree of function \code{copy}}
%         \label{fig:copy:SET}
% \end{figure}

%}}}
%{{{ Instrumentation

\section{Instrumentation} \label{sec:fmics_instr}

In our algorithm, the purpose of the instrumentation is to insert code
implementing an automaton into the analysed program. Nonetheless, the
semantics of the program being instrumented must not be changed. A result of
this phase is therefore a new program that still has the original
functionality but it also simultaneously updates instrumented state machines.
We show the process using the automaton $\SA(x)$ of Figure~\ref{fig:sm} and a
program consisting of two functions: \code{copy} of Figure~\ref{fig:copy} and
\code{foo} of Figure~\ref{fig:fmics_foo}.
\begin{figure}[htb]
  \begin{center}
  \begin{tabular}{c}
    \begin{lstlisting}[language=C]
char *buf1, *buf2;
int L1, L2;

void foo(char *src, int n) {
  copy(buf1, src, n, &L1);
  copy(buf2, src, n, &L2);
}
    \end{lstlisting}
  \end{tabular}
  \end{center}
  \caption{The function \code{foo} forms an analysed program together
    with the function \code{copy}.}
  \label{fig:fmics_foo}
\end{figure}
The function \code{foo} calls \code{copy} twice, first with the lock \code{L1}
and then with the lock \code{L2}. The locks protect writes into buffers
\code{buf1} and \code{buf2} respectively.  The function \code{foo} is a
so-called \emph{starting function}. It is a function where the symbolic
execution starts.

The instrumentation starts by recognizing code fragments which manipulate with
locks in the analysed program. More precisely, we look for all those code
fragments matching edge labels of the automaton $\SA(x)$ of
Figure~\ref{fig:sm}. The analysed program contains three such fragments, all
of them in the function \code{copy} (see Figure~\ref{fig:copy}): the call to
\code{lock} on line~\ref{line:copy_lock}, the call to \code{unlock} on
line~\ref{line:copy_unlock}, and the \code{return} statement on
line~\ref{line:copy_ret}.

Next we determine a set of all locks that are manipulated by the program.
From the recognized code fragments, we find out that the pointer variable
\code{L} in \code{copy} is the only program variable through which the
program manipulates with locks.  Using the pointer analysis, we obtain the set
$\{ \var{L1}, \var{L2} \}$ of all possible locks the program manipulates with.

We introduce a unique instance of the automaton $\SA(x)$ for each lock in the
set. More precisely, we define two integer variables \code{smL1} and
\code{smL2} to keep the current state of automata $\SA(\var{L1})$ and
$\SA(\var{L2})$, respectively. As we represent both automaton states and names
of transitions by integer constants, we need to define them too. So the
constants and the lock states are shown in Figure~\ref{fig:fmics_const}.
\begin{figure}[htb]
  \begin{center}
  \begin{tabular}{c}
    \begin{lstlisting}[name=fmics_impl]
const int smU  = 0;     // state U
const int smL  = 1;     // state L
const int smDU = 2;     // state DU
const int smDL = 3;     // state DL
const int smRL = 4;     // state RL

const int smLOCK   = 0; // transition lock(x)
const int smUNLOCK = 1; // transition unlock(x)
const int smRETURN = 2; // transition return

int smL1 = smU, smL2 = smU;
    \end{lstlisting}
  \end{tabular}
  \end{center}
  \caption{Constants and lock variables.}
  \label{fig:fmics_const}
\end{figure}
And we see that the lock states are set to the initial state of the automaton
there too.  Further, we need to specify a mapping from locks to their
automata. The mapping is basically a function (preferably with constant
complexity) from addresses of program objects (i.e.\ the locks) to addresses
of corresponding automata.  Figure~\ref{fig:fmics_getm} shows an
implementation of the function \code{smGetMachine} that maps addresses of
locks \code{L1} and \code{L2} to addresses of corresponding automata.
\begin{figure}[htb]
  \begin{center}
  \begin{tabular}{c}
    \begin{lstlisting}[name=fmics_impl]
int *smGetMachine(int *p) {
  if (p == &L1) return &smL1;
  if (p == &L2) return &smL2;
  return NULL;         // unreachable
}
    \end{lstlisting}
  \end{tabular}
  \end{center}
  \caption{An implementation of the automaton identification.}
  \label{fig:fmics_getm}
\end{figure}
We note that the implementation of \code{smGetMachine} would be more
complicated if automata were associated to dynamically allocated objects.
Although we have already instrumented the program with state machines, there
is no code manipulating them. For that purpose, we have the function
\code{smFire} implementing the automaton $\SA(x)$ in
Figure~\ref{fig:fmics_fire}.  Also keep in mind that the pointer argument
\code{SA} of the function points to an instrumented automaton, whose
transition has to be fired.

\begin{figure}[htb]
  \begin{center}
  \begin{tabular}{c}
    \begin{lstlisting}[name=fmics_impl]
void smFire(int *SA, int transition) {
  switch (*SA) {
  case smU:
    switch (transition) {
    case smLOCK:
      *SA = smL;
      break;
    case smUNLOCK:
      assert(false); // double unlock(*@\label{line:fmics_err_du}@*)
      break;
    default: break;
    }
  break;
  case smL:
    switch (transition) {
    case smLOCK:
      assert(false); // double lock(*@\label{line:fmics_err_dl}@*)
      break;
    case smUNLOCK:
      *SA = smU;
      break;
    case smRETURN:
      assert(false); // return in locked(*@\label{line:fmics_err_l}@*)
      break;
    default: break;
    }
    break;
  default: break;
  }
}
    \end{lstlisting}
  \end{tabular}
  \end{center}
  \caption{The implementation of the automaton transitions.}
  \label{fig:fmics_fire}
\end{figure}

It remains to instrument the recognized code fragments in the original
program. For each fragment we know its related transition of the automaton and
we also know what objects the fragment manipulates with (if any). Therefore,
we first retrieve an address of automaton related to manipulated objects (if
any) by using the function \code{smGetMachine} and then we fire the transition
by calling the function \code{smFire}. The instrumented version of the
original program consists of constants with lock variables
(Figure~\ref{fig:fmics_const}), the functions \code{smGetMachine}
(Figure~\ref{fig:fmics_getm}) and \code{smFire} (Figure~\ref{fig:fmics_fire}),
and the instrumented version of the original functions \code{foo} and
\code{copy} given in Figure~\ref{fig:copy_instrumented}. The instrumented
lines are highlighted by~\code{*}.  Note that in our example, the instrumented
state variables \code{smL1} and \code{smL2} directly correspond to the program
locks \code{L1} and \code{L2} respectively. In general, nonetheless, states of
a state machine of a more complex property need not necessarily correspond to
values of a particular program variable. Therefore, we apply this general
approach to our example too.

\begin{figure}[tb]
    \begin{center}
    \begin{tabular}{l}
    \begin{lstlisting}[language=C]
  char *buf1, *buf2;
  int L1, L2;

  char *copy(char *dst, char *src, int n, int *L) {
    int i, len;
    len = 0;
    if (src != NULL && dst != NULL) {
      len = n;
*     smFire(smGetMachine(L), smLOCK);
      lock(L);
    }
    i = 0;
    while (i < len) {
      dst[i] = src[i];
      i++;
    }
    if (len > 0) {
*     smFire(smGetMachine(L), smUNLOCK);
      unlock(L);
    }
*   smFire(smGetMachine(L), smRETURN);
    return dst;
  }

  void foo(char *src, int n) {
    copy(buf1, src, n, &L1);
    copy(buf2, src, n, &L2);
  }
    \end{lstlisting}
    \end{tabular}
    \end{center}
    \caption{Functions \code{foo} and \code{copy} instrumented by calls
      of \code{smFire} function. The instrumented code is marked by \code{*}.}
    \label{fig:copy_instrumented}
\end{figure}

Note that we instrumented the program of the figures somewhat manually. In
practice, we need this process automated somehow. One solution is to ask a
user to instrument header files and a preprocessor distributes the transition
functions automatically. This approach was successfully employed for checking
of the Linux kernel locking policy. We manually instrumented macros for
locking like \code{spin_lock} and \code{spin_unlock}. And the preprocessor
during the build process expanded these macros in every place where locking is
used and put there the function \code{smFire}. But as one can imagine, this is
only possible for projects and properties where such macros exist. A use case
where this does not work are properties regarding pointer manipulation errors.
There, we need to fire a transition by most pointer manipulations, but barely
someone will put some macros around them. In that case, one can use
instrumenting tools like the one from the \textsc{Ldv} toolkit~\cite{ldv}.
They perform the model checking of various properties in the kernel and have
the same use case of instrumentation as we do.

%}}}
%{{{ Slicing

\section{Slicing} \label{sec:fmics_slicing}

Let us have a look at the instrumented program in
Figure~\ref{fig:copy_instrumented}. We can easily observe, that the main part
of the function \code{copy}, the loop copying the characters, does not
affect states of the instrumented automata. The symbolic execution of such
code is known to be very expensive, moreover in this case it is yet unneeded.
Therefore, we use the slicing technique from Weiser~\cite{slicing} to
eliminate such code from the instrumented program.

\begin{figure}[bt]
    \begin{center}
    \begin{tabular}{l}
    \begin{lstlisting}[language=C]
char *buf1, *buf2;
int L1, L2;

char *copy(char *dst, char *src, int n, int *L) {
  int len;
  len = 0;
  if (src != NULL && dst != NULL) {
    len = n;
    smFire(smGetMachine(L), smLOCK);
  }
  if (len > 0) {(*@\label{line:fmics_len_cond}@*)
    smFire(smGetMachine(L), smUNLOCK);
  }
  smFire(smGetMachine(L), smRETURN);(*@\label{line:fmics_2nd_fire}@*)
  return dst;
}

void foo(char *src, int n) {
   copy(buf1, src, n, &L1);
   copy(buf2, src, n, &L2);
}
    \end{lstlisting}
    \end{tabular}
    \end{center}
    \caption{Functions \code{foo} and \code{copy} after slicing.}
    \label{fig:copy_sliced}
\end{figure}

Our analysis is interested only in states of the instrumented automata,
especially in locations corresponding to errors.  Hence, the slicing criterion
is a pair of a location preceding an \code{assert} statement in the
\code{smFire} function and the set of all variables representing current
states of the corresponding automata. The slicing criteria then comprises all
such pairs.

In the instrumented program consisting of Figures~\ref{fig:fmics_const},
\ref{fig:fmics_getm}, \ref{fig:fmics_fire}, and~\ref{fig:copy_instrumented},
we want to preserve variables \code{smL1} and \code{smL2}. We put slicing
criteria to the lines of the code detecting transitions of automata into error
states. In other words, the slicing criteria for our running example are pairs
(\ref{line:fmics_err_du}, \{\code{smL1},\code{smL2}\}),
(\ref{line:fmics_err_dl}, \{\code{smL1},\code{smL2}\}), and
(\ref{line:fmics_err_l}, \{\code{smL1},\code{smL2}\}), where the numbers refer
to lines in the code of Figure~\ref{fig:fmics_fire}.  The result of the
slicing procedure is presented in Figures~\ref{fig:fmics_const},
\ref{fig:fmics_getm}, \ref{fig:fmics_fire}, and~\ref{fig:copy_sliced} (the
code in the first three figures shall not be changed by the slicing). Note
that the sliced code contains neither the \code{while}-loop nor the
\code{lock()} and \code{unlock()} commands.

It is important to note that some slicing techniques, including the one
of Weiser~\cite{slicing} which we use, do not consider inputs for which the
original program does not halt.
% The problem arises when the sliced code contains an error, such that all
% executions exhibiting the error do not reach the error in the original
% code as they loop in some code that sliced.
As a result, an input causing an infinite run of the original program can
induce a finite run in the sliced program. Moreover, the finite run can
contain locations not visited by the infinite run.
% Therefore, there is no way to guarantee that a sliced program would fail
% to halt whenever the original program fails to halt.
This is the only principal source of potential false positives in our
technique.

%}}}
%{{{ Symbolic Execution

\section{Symbolic Execution} \label{sec:fmics_symexe}

The symbolic execution is the final phase of our technique. We symbolically
execute the sliced program from the entry location of the starting function.
The symbolic execution explores real program paths. Therefore, if it reaches
some of the assertions inside function \code{smFire}, then we have found a
bug.

Our running example nicely illustrates the crucial role of slicing to
feasibility of the symbolic execution. Let us first consider a symbolic
execution of the original program. It starts at the entry location of the
function \code{foo}. The execution eventually reaches the function
\code{copy}. Note that value of the parameter \code{n} is symbolic. Therefore,
the symbolic execution will fork into two executions each time we reach
line~\ref{line:copy_loop} of Figure~\ref{fig:copy}. One of the executions
skips the loop at lines~\ref{line:copy_loop}--\ref{line:copy_loopE},
while the other enters it. If we assume that the type of \code{n} is a 32-bit
integer, then the symbolic execution of one call of \code{copy} explores more
then $2^{31}$ real paths.

\begin{figure}[!tb]
    \begin{center}
    \tikzstyle{loc} = [circle, thick, draw, minimum size=6mm, inner sep=0pt]
    \tikzstyle{pre} = [<-, shorten <=1pt, >=stealth', semithick]
    \tikzstyle{post} = [pre, ->]
    \begin{tikzpicture}[node distance=4mm and 5mm, font=\footnotesize]
        \node [loc] (0) {$19$};
        \node [loc] (1) [below=of 0] {$6$}
            edge [pre] node {} (0);
        \node [loc] (2) [below=of 1] {$7$}
            edge [pre] node {} (1);

        % ->false branch
        \node [loc] (3) [below left=4mm and 14mm of 2] {$11$}
            edge [pre] node [above left, xshift=1mm] {\code{false}} (2);
        \node [loc] (4) [below=of 3] {$14$}
            edge [pre] node {} (3);
        \node [loc] (5) [below=of 4] {$15$}
            edge [pre] node {} (4);
        \node [loc] (20) [below=of 5] {$20$}
            edge [pre] node {} (5);
        \node [loc] (21) [below=of 20] {$6$}
            edge [pre] node {} (20);
        \node [loc] (22) [below=of 21] {$7$}
            edge [pre] node {} (21);
        \node [loc] (23) [below right=of 22] {$8$}
            edge [pre] node [above right, xshift=-1mm] {\code{true}} (22);
        \node [loc] (24) [below=of 23] {$9$}
            edge [pre] node {} (23);
        \node [loc] (25) [below=of 24] {$11$}
            edge [pre] node {} (24);
        \node [loc] (26) [below left=of 22] {$11$}
            edge [pre] node [above left, xshift=1mm] {\code{false}} (22);
        \node [loc] (27) [below=of 26] {$14$}
            edge [pre] node {} (26);
        \node [loc] (28) [below=of 27] {$15$}
            edge [pre] node {} (27);
        \node [loc] (29) [below left=of 25] {$14$}
            edge [pre] node [above left, xshift=1mm] {\code{false}} (25);
        \node [] (BUG) [below=1mm of 29] {\emph{error}};
        \node [loc] (30) [below right=of 25] {$12$}
            edge [pre] node [above right, xshift=-1mm] {\code{true}} (25);
        \node [loc] (31) [below=of 30] {$14$}
            edge [pre] node {} (30);
        \node [loc] (32) [below=of 31] {$15$}
            edge [pre] node {} (31);

        % ->true branch
        \node [loc] (7) [below right=4mm and 14mm of 2] {$8$}
            edge [pre] node [above right, xshift=-1mm] {\code{true}} (2);
        \node [loc] (8) [below=of 7] {$9$}
            edge [pre] node {} (7);
        \node [loc] (9) [below=of 8] {$11$}
            edge [pre] node {} (8);

        % ->true->false branch
        \node [loc] (11) [below left=of 9] {$14$}
            edge [pre] node [above left, xshift=1mm] {\code{false}} (9);
        \node [] (BUG) [below=1mm of 11] {\emph{error}};

        % ->true->true branch
        \node [loc] (13) [below right=of 9] {$12$}
            edge [pre] node [above right, xshift=-1mm] {\code{true}} (9);
        \node [loc] (14) [below=of 13] {$14$}
            edge [pre] node {} (13);
        \node [loc] (15) [below=of 14] {$15$}
            edge [pre] node {} (14);

        \node [loc] (40) [below=of 15] {$20$}
            edge [pre] node {} (15);
        \node [loc] (41) [below=of 40] {$6$}
            edge [pre] node {} (40);
        \node [loc] (42) [below=of 41] {$7$}
            edge [pre] node {} (41);
        \node [loc] (43) [below right=of 42] {$8$}
            edge [pre] node [above right, xshift=-1mm] {\code{true}} (42);
        \node [loc] (44) [below=of 43] {$9$}
            edge [pre] node {} (43);
        \node [loc] (45) [below=of 44] {$11$}
            edge [pre] node {} (44);
        \node [loc] (46) [below right=of 45] {$12$}
            edge [pre] node [above right, xshift=-1mm] {\code{true}} (45);
        \node [loc] (47) [below=of 46] {$14$}
            edge [pre] node {} (46);
        \node [loc] (48) [below=of 47] {$15$}
            edge [pre] node {} (47);
        \node [loc] (49) [below left=of 42] {$11$}
            edge [pre] node [above left, xshift=1mm] {\code{false}} (42);
        \node [loc] (50) [below=of 49] {$14$}
            edge [pre] node {} (49);
        \node [loc] (51) [below=of 50] {$15$}
            edge [pre] node {} (50);

%        \node [] (X1) [above of=1,yshift=2mm] {}
%            edge [dotted,thick,draw] node {} (1);
    \end{tikzpicture}
    \end{center}
    \caption{A symbolic execution tree of the sliced program of
      Figure~\ref{fig:copy_sliced}.}
        \label{fig:copy_sliced:SET}
\end{figure}

By contrast, the sliced program does not contain the loop, which generated
the huge number of real paths. Therefore, a number of real paths explored by
the symbolic execution is exactly \emph{six}. Figure~\ref{fig:copy_sliced:SET}
shows the symbolic execution tree of the sliced program of
Figure~\ref{fig:copy_sliced}. We left out nodes corresponding to lines in
called functions \code{smGetMachine} and \code{smFire}. Note that
although the parameter \code{n} has a symbolic value, it can only affect
the branching on line~\ref{line:fmics_len_cond}. Moreover, the parameter
\code{L} always has a concrete value. Therefore, we do not fork the symbolic
execution at branchings inside functions \code{smGetMachine} and
\code{smFire}. Two of the explored paths are marked with the label
\emph{error}.  These paths reach the second assertion in function
\code{smFire} (see Figure~\ref{fig:fmics_fire}) called from
line~\ref{line:fmics_2nd_fire} of the sliced program. In other words, the
paths are witnesses that we can leave the function \code{copy} in a locked
state. The remaining explored paths of Figure~\ref{fig:copy_sliced:SET} miss
the assertions in the function \code{smFire}. It means that the original
program contains only one locking error, namely \emph{return in locked state}.

It might be the case for some program and checked property, that the sliced
code still contains loops. Then the subsequent symbolic execution can be very
costly due to the well-known path explosion problem. Fortunately, there have
been done some work tackling the problem~\cite{G07, grammar3, GNRT10, SPmCS09,
XGM08, cse_atva}.

% analysis.

%}}}
%{{{ Implementation

\section{Implementation} \label{sec:fmics_impl}
To verify the applicability of the presented technique, we have developed an
experimental implementation. We call the experimental tool \symbio. \symbio
works with programs written in C and, for the sake of simplicity, it detects
only locking errors described by an automaton very similar to $\SA(x)$ of
Figure~\ref{fig:sm}. The instances of the automaton are associated with
arguments of \code{lock} and \code{unlock} function calls. Note that the
technique currently works only for the cases where a lock is instantiated only
once during the run of the symbolic executor, which is the most frequent case.
Nonetheless we plan to add a support even for the rest. The main part of our
implementation is written in three modules for the \llvm
framework~\cite{llvm}, namely \code{Prepare}, \code{Slicer}, and
\code{Kleerer}. The framework provides us with a C compiler
\clang~\cite{clang}. We also use an existing symbolic executor for \llvm
called \klee~\cite{klee}. The whole pipeline is depicted in
Figure~\ref{fig:fmics_pipeline}

\begin{figure}[tb]
    \begin{center}
    \begin{tikzpicture}[node distance=6mm, font=\sffamily]
      \node (S) {};
      \node [tl_state, right=5mm of S] (C) {C Sources}
        edge[tl_edge] node {} (S);
      \node [tl_phase, right=of C, align=center] (PREP) {
	  C Preprocessor/ \\
	  Instrumentor
	}
        edge[tl_edge] node {} (C);
      \node [tl_state, right=of PREP, align=center] (PREPS) {
	  Instrumented \\
	  and preprocessed \\
	  C code
	}
        edge[tl_edge] node {} (PREP);
      \node [tl_phase, below=6mm of PREPS] (CLANG) {\clang}
        edge[tl_edge] node {} (PREPS);
      \node [tl_state, left=of CLANG] (LLVM) {Llvm}
        edge[tl_edge] node {} (CLANG);
      \node [tl_phase, left=of LLVM] (PREPARE) {\code{Prepare}}
        edge[tl_edge] node {} (LLVM);
      \node [tl_state, left=of PREPARE, align=center] (LLVMPREP) {
	  Llvm with \\
	  globals etc.
	}
        edge[tl_edge] node {} (PREPARE);
      \node [tl_phase, left=of LLVMPREP] (SLICER) {\code{Slicer}}
        edge[tl_edge] node {} (LLVMPREP);
      \node [tl_state, below=8mm of SLICER, align=center] (SLICED) {Sliced \\ Llvm}
        edge[tl_edge] node {} (SLICER);
      \node [tl_phase, right=of SLICED] (KLEERER) {\code{Kleerer}}
        edge[tl_edge] node {} (SLICED);
      \node [tl_state, right=of KLEERER, align=center] (KLEERED) {
	  Llvm with \\
	  \code{main}
	}
        edge[tl_edge] node {} (KLEERER);
      \node [tl_phase, right=of KLEERED] (KLEE) {\klee}
        edge[tl_edge] node {} (KLEERED);
      \draw [tl_edge] ([yshift=-2mm] KLEE.west) -- ([yshift=-2mm]KLEERED.east);
      \draw [tl_edge] ([yshift=2mm] KLEE.west) -- ([yshift=2mm]KLEERED.east);
      \node [tl_state, right=of KLEE] (RESULTS) {Results}
        edge[tl_edge] node {} (KLEE);

      \node [inner sep=7pt, fit=(PREP)] (PREPBB) {};
      \node [inner sep=7pt, fit=(PREPARE)] (PREPAREBB) {};
      \path[draw, densely dashed]
	(PREPBB.north west) -- (PREPBB.south west) --
	(PREPAREBB.north west) -- (PREPAREBB.south west) --
	(PREPAREBB.south east) -- (PREPAREBB.north east) --
	(PREPBB.south east) -- (PREPBB.north east) -- (PREPBB.north west);
      \node [above=-0.5mm of PREPARE, align=center, xshift=1mm] (2PHASE_DESC) {
	  \footnotesize{Instrumentation} \\[-1.5mm]
	  \footnotesize{(2 phases)}
	};
    \end{tikzpicture}
    \end{center}
    \caption{The pipeline of \symbio.}
    \label{fig:fmics_pipeline}
\end{figure}

\subsection{Prepare}
The instrumentation of a given program proceeds in two steps as can be seen in
Figure~\ref{fig:fmics_pipeline}. Using a standard C preprocessor, the original
program is instrumented with function calls to \code{smFire} located just
above statements changing states of automata.  The program is then translated
by \clang into the \llvm bytecode~\cite{llvm}. Since \klee requires \llvm code
to be compiled without optimizations, \clang is set to generate such code.
Then, the final stage of instrumentation is performed on the \llvm code by
the module \code{Prepare}. This phase performs several tasks:
\begin{itemize}
\item It adds missing global variables and makes them symbolic.
\item It removes bodies of many predefined functions, because they contain
assembler. Instead, it links in pure \llvm implementations.
\item It removes calls to functions which have no side effect. These include
printing to the output (\code{printf}).
\item It finds assembler patterns and replaces them by pure \llvm code.
\item Finally, it establishes the set of starting functions similarly to \stan
(Chapter~\ref{sec:stanse}): roots in the call graph.
\end{itemize}

\subsection{Slicer}
The module \code{Slicer} implements a variant of the inter-procedural
slicing algorithm by Weiser~\cite{slicing} as described in
Section~\ref{sec:slicing}. Note that Weiser's algorithm poses a requirement on
control flow graphs which are to be sliced. They need to be \emph{hammock
graphs}, i.e.\ graphs with a single entry and single exit. As \clang does not
necessarily generate such control flow graphs, we convert the graphs into this
structure first. There is a submodule called \code{CreateHammockCFG} for this
purpose.

\subsubsection{LLVM Pointer Analysis}
To guarantee correctness and to improve performance of the slicing, the
algorithm employs the pointer analysis by Andersen~\cite{andersen}. For the
better understanding and possible replacement, we convert the \llvm code first
into a sequence of statements like \code{VAR(x)=ALLOC}, \code{VAR(x)=VAR(y)},
and \code{DEREF(VAR(x))=REF(VAR(y))} according to the \llvm semantics.  Only
these are then passed to the pointer analysis. The implementation of the
analysis then computes a fixed point over the statements using the Andersen's
rules. Finally, pointers are in these terms \llvm variables and their pointer
sets contain global variables and \llvm instructions where an allocation
happen.

Initially, we had only a field-insensitive analysis, but this caused too much
code to be sliced away. As an example, it insisted on uttering that the two
statements \code{a.x = 0} and \code{a.y = 0} store to the same memory
location. So the pointer analysis allowed the slicer to remove the first
assignment, while it still perfectly can affect user's slicing criteria. Hence
pointers were changed to be pairs of a \llvm variable and byte offset. The
same holds for members of pointer sets. One may wonder why we need pairs for
pointers, why not only for a pointed memory. But consider a structure in C.
Each member can be a pointer and that can point somewhere. We can now have a
separate information for each of the pointer in that structures. The extension
is straightforward in the scope of \llvm. It has a special instruction for
taking an offset of a memory location: \code{GetElementPtr}. So we added
another rule to the pointer analysis framework \code{VAR(x)=GEP(y)}. Finally,
the pointer analysis implementation extracts the offset from the instruction's
parameters and adds it to the offset member of the pointed pair.

\subsection{Kleerer}
The module \code{Kleerer} performs a final processing of the sliced bytecode
before it is passed to \klee. In particular, the module adds to the bytecode a
function \code{main} that calls a starting function.  The \code{main}
function also allocates a symbolic memory for each parameter of the starting
function. Size of the allocated memory is determined by the parameter type.
More problematic is a pointer as an argument of the starting function. We have
no way to find out whether the pointer points to a single value in memory or
whole array. To overcome the problem, we assume there are at most 4192 members
in the array. So for example, 4 bytes are allocated for an integer and 16768
bytes for an integer pointer. The idea behind the array allocations is taken
over from the paper introducing \ucklee~\cite{ucklee}.

We still had to extend the approach by two means. First, we enumerate some
known structures which are never passed as arrays. An example of such
structure is a PCI device descriptor stored in the structure \code{pci_dev}.
We just pass a memory of the structure size.  Second, we pass to a starting
function a pointer to the middle of the allocated memory. This is because
functions may dereference memory at a negative index. This usually happens
when some callback in a driver is called and it has as a parameter some kernel
structure. Such structure is usually a static member of another structure
owned by a driver. Let us illustrate that on an example taken from the Linux
kernel 3.11, file \code{drivers/pci/hotplug/pciehp_ctrl.c}, i.e.\ a PCI
Express hot-plug driver.
\begin{figure}[htb]
    \begin{center}
    \begin{tabular}{c}
    \begin{lstlisting}
struct power_work_info {(*@\label{line:fmics_neg_struct}@*)
  struct slot *p_slot;
  struct work_struct work;(*@\label{line:fmics_neg_kernel}@*)
};

void pciehp_power_thread(struct work_struct *work) {
  struct power_work_info *info =
	  container_of(work, struct power_work_info, work);(*@\label{line:fmics_neg_neg}@*)
  ...
    \end{lstlisting}
    \end{tabular}
    \end{center}
    \caption{Demonstration of the negative pointer index in a starting
    function.}
    \label{fig:fmics_neg}
\end{figure}
There is a declaration of the driver structure \code{power_work_info} on
line~\ref{line:fmics_neg_struct}. The structure contains also the kernel
internal structure \code{work_struct} on line~\ref{line:fmics_neg_kernel}.
Finally, when the callback \code{pciehp_power_thread} is called with a pointer
to \code{work_struct}, the callback uses \code{container_of} to extract back a
pointer to \code{power_work_info}. Obviously, the pointer to
\code{power_work_info} is smaller than to \code{work_struct}.

\subsection{Symbolic Execution}
After we obtain \llvm code with the \code{main} function, we can start the
symbolic execution. Note that in principle, there can be more than one
starting function, i.e.\ more input files for the symbolic execution. But
\klee needs to perform the symbolic execution only on one starting function.
On the one that was used in the analyzer that generated a report which we are
trying to verify.

The symbolic execution need not finish if the analyzed code is still too
complex. But ideally, when if finishes, it can do so in several ways. First,
\klee can tell us there is no error or there is an error. We have a conclusion
in that case.  But \klee also sometimes crashes, or finishes the execution
with a \emph{memory error}.  In either case, we cannot conclude anything about
the input file. The memory errors are caused by too low limits for the
allocated memory and the code touched a memory out of bounds while was being
executed.  To remedy this inconvenience, we plan to implement the same
on-demand memory handling as \ucklee does.

\section{Experimental Results}\label{sec:fmics_experiments}

\begin{table}[tb]
\begin{center}
\setlength{\tabcolsep}{1.2pt}
\renewcommand*\arraystretch{1.2}
%\scriptsize
\begin{tabular}{|l|c|c|c|c|c|c|c|c|} \hline
  {\textbf{File}}&
  \multicolumn{5}{c|}{\textbf{Running Time} (s)} & 
  \multirow{2}{*}{\textbf{Sliced}} &
  \multirow{2}{*}{\textbf{Res.}} & \textbf{Factual} \\ \cline{2-6}
  {\textbf{Function}} & \textbf{Co.} & \textbf{Inst.} &
   \textbf{Slic.} & \textbf{SE} & \textbf{Total} & &
 & \textbf{State} \\ \hline \hline

%  \begin{minipage}{3cm}
%    \smallskip\scriptsize{\code{fs/jfs/super.c}} \\[-1pt]
    \code{jfs_quota_write}%\smallskip
%  \end{minipage}
  & 1.25 & 0.18 & 0.15 & 5.09 & 6.67 & 
  \twolines{369/119}{67.8\%} & E & E \\ \hline

  \begin{minipage}{3cm}
%    \smallskip\scriptsize{\code{drivers/net/qlge/qlge_main.c}} \\[-1pt]
    \code{qlge_set_mac_address}\smallskip
  \end{minipage}
  & 2.70 & 0.72 & 26.75 & 13.28 & 43.45 &
  \twolines{1333/447}{66.5\%} & E & E \\ \hline

  \begin{minipage}{3cm}
%    \smallskip\scriptsize{\code{drivers/hid/hidraw.c}} \\[-1pt]
    \code{hidraw_read}\smallskip
  \end{minipage}
  & 1.06 & 0.18 & 0.14 & \multicolumn{2}{c|}{Timeout} & 
  \twolines{666/220}{67.0\%} & TO & E \\ \hline

  \begin{minipage}{3cm}
%    \smallskip\scriptsize{\code{drivers/net/ns83820.c}} \\[-1pt]
    \code{queue_refill}\smallskip
  \end{minipage}
  & 1.76 & 0.29 & 1.72 & 0.62 & 4.39 & 
  \twolines{1212/329}{72.9\%} & FP & FP \\[-2pt] \hline

  \begin{minipage}{3.6cm}
%    \smallskip\scriptsize{\code{drivers/usb/misc/}} \\[-3pt]
%    \hspace*{2ex}\scriptsize{\code{sisusbvga/sisusb_con.c}}\\[-1pt]
    \code{sisusbcon_set_palette}\smallskip
  \end{minipage}
  & 1.50 & 0.24 & 0.27 & 17.19 & 19.20 & 
  \twolines{2936/705}{76.0\%} & FP & FP \\\hline

  \begin{minipage}{3cm}
%    \smallskip\scriptsize{\code{fs/jffs2/nodemgmt.c}} \\[-1pt]
    \code{jffs2_reserve_space}\smallskip
  \end{minipage}
  & 1.04 & 0.18 & 0.22 & \multicolumn{2}{c|}{Timeout} & 
  \twolines{677/360}{46.8\%} & TO & FP \\ \hline

  \begin{minipage}{3cm}
%    \smallskip\scriptsize{\code{kernel/kprobes.c}} \\[-1pt]
    \code{pre_handler_kretprobe}\smallskip
  \end{minipage}
  & 0.32 & 0.09 & 0.51 & 2.43 & 3.35 & 
  \twolines{202/68}{66.3\%} & ME & FP \\ \hline

\end{tabular}
\end{center}
\caption{Experimental results. The table presents running time of
  preprocessing and compilation (\textbf{Co.}), instrumentation including
  points-to analysis (\textbf{Inst.}), slicing (\textbf{Slic.}), symbolic
  execution (\textbf{SE}), and the total running time. The column
  \textbf{Sliced} presents the ratio of instructions sliced away from the
  instrumented \textsc{Llvm} code and the exact number of instructions
  before/after slicing. The column \textbf{Res.} specifies the
  result of our tool: E means that the tool found a real error, FP means
  that the analysis finished without error found (i.e.~the original error
  report is a false positive), TO that the symbolic execution did not finish in
  time and ME denotes an occurrence of memory error. The last column specifies
  the factual state of the error report.}
\label{tab:results_fmics}
\end{table}

% fs/jfs/super.c jfs_quota_write 369 / 119 (67.75% is gone)
% no slicing: timeout
% drivers/net/qlge/qlge_main.c qlge_set_mac_address 1333 / 447 (66.47% is gone)
% no slicing: timeout
% drivers/net/ns83820.c queue_refill 1212 / 329 (72.85% is gone)
% no slicing: timeout
% drivers/usb/misc/sisusbvga/sisusb_con.c sisusbcon_set_palette 2936 / 705 (75.99% is gone)
% no slicing: 20.64 s
% fs/jffs2/nodemgmt.c jffs2_reserve_space (to po 5m40s) 677 / 360 (46.82% is gone)
% no slicing: timeout
% drivers/hid/hidraw.c hidraw_read (to po 10m) 666 / 220 (66.97% is gone)
% no slicing: timeout
% kernel/kprobes.c pre_handler_kretprobe (mem err) 202 / 68 (66.34% is gone)
% no slicing: timeout
% klee -max-stp-time=10 -max-time=300
We have performed our experiments on several functions of the Linux kernel
2.6.28, where the static analyzer \textsc{Stanse} reported some error. More
precisely, \textsc{Stanse} reported an error trace starting in these
functions. We consulted the errors with kernel developers to sort out which
are false positives and which are real errors. All the selected functions
(and all functions transitively called from them) contain no assembler (in
some cases, it has been replaced by an equivalent C code) and no external
function calls after slicing.

We ran our experimental tool on these functions. All tests were performed on
a machine with an Intel E6850 dual-core processor at 3\,GHz and 6\,GiB of
memory, running Linux. We specified \klee parameters to time out after
10 seconds spent in an SMT solver and after 300 seconds of an overall running
time. Increasing these times brings no real effect in our environment. We do
not pass optimize option for \klee because it causes \klee to
crash for most of the input.
% The optimization might
% further speed the analysis up. The crashes are yet to be investigated.

% At this early stage of the tool, the evaluation is not performed on the whole
% kernel. It is because it still does not scale that much. We plan to extend
% the tool to be capable of handling whole kernel. We limited our testing to a
% set of files. We aimed the technique at the functions where abstract
% interpretation tools reported a trace. Be it a real error or false positive.
% We consulted the files with upstream developers to sort out which are in each
% category.

Table~\ref{tab:results_fmics} presents results of our tool on selected
functions.  The table shows compilation, instrumentation, slicing, symbolic
execution, and the overall running time. Further, the table presents the ratio
of instructions that were sliced away from the instrumented \textsc{Llvm}
code.  The last two columns specify the results of our analysis and the real
state confirmed by kernel developers. The table clearly shows that the
bottleneck of our technique is the symbolic execution. However, if we did not
slice the code, the only function completely executed in time would be
\code{sisusbcon\_set\_palette}, computed in 20.64s.
% , i.e. slower by 17 per cent. As a demonstration, slicer reduces 2936
% \textsc{Llvm} instructions to 705 there.

Although the results have no statistical significance, it is clear that the
technique can in principle classify error reports produced by other tools like
\textsc{Stanse}. If our technique reports an error, it is a real one. If it
finishes the analysis without any error detected, the original error report is
a false positive. The analysis may also not finish in a given time, which is
usually caused by loops in the sliced code. Finally, it may report a memory
error mentioned above.


% In most cases with our experiments, the time spent by compilation,
% instrumentation and slicing may be silently ignored with respect to the SE.
% In the kernel, there are few exception like \texttt{qlge\_set\_mac\_address}
% presented here.  Slicing in this case is a bit expensive due to circular
% inter-procedural dependencies of the slicing criteria. So that many states are
% transferred over calls until the fixed point is reached. In any case, the time
% of slicing was not more than a minute.

% Other than that, the vast majority of the analysis time is spent with the SE.
% In some cases, the time limit is hit. In both functions in
% Table~\ref{tab:results} it it because of unbounded loops which depend on
% locks.

% Table~\ref{tab:results} also sketches an effect of the slicing. It can be seen
% that it is effective as it allows, in average, 66 per cent of the code to be
% removed.

%}}}
%{{{ Related Work

%\subsection{Related Work} \label{sec:related}
%
%% A description of program properties in \textsc{Metal} language and
%% meta-level compilation is discussed in~\cite{CEH02,metacomp,ECC00,metal}. 
%
%There are many tools checking properties described by finite state
%machines. They produce both kinds of reports, real error as well as false
%positives. The technique of \textsc{xgcc} presented
%in~\cite{CEH02,metacomp,ECC00,metal} found a thousands of bugs in real system
%code. It provides a language \textsc{Metal} for easy description of properties
%to be checked.  \textsc{xgcc} suffers from false positives despite usage of
%false positive suppression algorithms like killing
%variables and expressions, synonyms, false path pruning, and others. Besides
%the suppression algorithms, bug-reports from the tool are further ranked
%according to their probability of being real errors. There are generic and
%statistical ranking algorithms ordering bug-reports. An extension introduced
%in~\cite{ECH01} provides an automatic inference of some temporal properties
%based on statistical analysis of assumed programmer's beliefs. The
%\textsc{ESP}~\cite{esp} technique uses a similar language to
%\textsc{Metal} for properties description. It implements an interprocedural
%dataflow algorithm based on~\cite{RHS95} for error detection and an abstract
%simulation pruning algorithm for false positives
%suppression. \textsc{Stanse}~\cite{stanse}, a static analysis tool also uses
%state machines for description of checked program properties. The
%description is based on parametrised abstract syntax trees.  Finally,
%\textsc{CEGAR}~\cite{cegar} based tools like \textsc{SLAM}~\cite{slam_decade},
%\textsc{SDV}~\cite{slam2}, or \textsc{Blast}~\cite{blast}, do not produce
%false positives, in theory. However, to achieve an appropriate efficiency
%and scalability for a practical use, the implementation of the
%\textsc{CEGAR} loop is typically unsound.
%%Although this
%%tool found hundreds of real bugs in the Linux kernel, it suffers from a high
%%false positive rate since its false positive suppression algorithms are
%%very limited.
%
%Program analysis tools based on symbolic execution~\cite{se} mainly discover
%low-level bugs like division by zero, illegal memory access, assertion
%failure etc. These tools typically do not have problems with false
%positives, but they have problems with scalability to large programs. There
%has been developed a lot of techniques improving the scalability to programs
%used in practice.

%There are also hybrid techniques combining symbolic execution with a
%complementary static
%analysis~\cite{GNRT10,NRTT09}. % Beckmanetal08,GMR09,Gulavanietal06,
%Symbolic execution can be accelerated by a compositional approach based on
%function summaries~\cite{AGT08,G07}. Another approach to effective symbolic
%execution introduced in~\cite{BCE08,klee,Cadar08} is based on recording of
%already seen behavior and pruning its repetition.
%% The followng techniques
%%focus on reaching a specific program location. Fitnex~\cite{XTHS09}, a
%%search strategy implemented in \textsc{Pex}~\cite{TdH08}, guides a path
%%exploration to a particular target location using fitness function. The
%%function measures how close an already discovered feasible path is to the
%%target. The LESE~\cite{SPmCS09} approach introduces symbolic variables for
%%the number of times each loop was executed. The symbolic variables are
%%linked with features of a known grammar generating inputs. Using these
%%links, the grammar can control the numbers of loop iterations performed on a
%%generated input. A technique presented in~\cite{GL11} analyses loops
%%on-the-fly, i.e.~during simultaneous concrete and symbolic executions of a
%%program for a concrete input. The loop analysis infers variables that are
%%modified by a constant value in each loop iteration. These variables are
%%used to build loop summaries expressed in a form of pre and post
%%conditions. An algorithm in~\cite{ST11} constructs a nontrivial necessary
%%condition on input values to drive the program execution to a given
%%location. A technique presented in~\cite{OT11} introduces a pair of counters
%%for two different paths around loop for each recurrent variable. Each
%%counter keeps an information about the number of iterations around one path
%%since the last iteration around the other one. Finally, 
%There is an orthogonal line of research which tries to improve the symbolic
%execution for programs with some special types of inputs. Some techniques
%deal with programs manipulating strings~\cite{BTV09,XGM08}, and some other
%techniques reduce input space using a given input
%grammar~\cite{grammar3,SPmCS09}.
%
%The interprocedural static slicing was introduced by Weiser~\cite{slicing}. But
%nowadays, there are many different approaches to program slicing. They are
%surveyed by several authors~\cite{BG,L01,T95}. Applications of slicing include
%program debugging, reverse engineering and regression testing~\cite{GHS92}.

%}}}
%{{{ Future Work

\section{Future Work}\label{sec:fmics_future}
Our future work has basically three independent directions.

First, we plan to run our tool to classify all lock-related error reports
produced by \textsc{Stanse} on the Linux kernel. The results should provide
a better image of practical applicability of the technique. To get a
relevant data, we should solve some practical issues like a correct
detection of starting functions, automatic replacement of assembler,
treatment of external function calls, etc. We should also implement an
on-demand memory allocation to \klee as discussed in
Section~\ref{sec:fmics_impl} or use a different executor.

The second direction is to adopt or design some convenient way for
specification of arbitrary state machines. It may be a dedicated language
similar to \textsc{Metal}~\cite{metal}. Then we plan to implement an
instrumentation treating these state machines. In particular, the
instrumentation should correctly handle state machines associated with
dynamically allocated objects.

Finally, we would also like to examine performance of our technique as a
stand-alone error-detection tool. To this point, we have to use a symbolic
executor aiming for maximal code coverage. In particular, such an executor
has to suppress execution paths that differ from explored paths only in
number of loop iterations. Unfortunately, we do not know about any publicly
available symbolic executor of this kind. However, it seems that
\ucklee~\cite{ucklee} (which is not public as of now) has been
designed for a similar purpose.

%}}}
%{{{ Conclusion

\section{Conclusion}\label{sec:fmics_conclusion}
We have presented a novel technique combining three standard methods
(instrumentation, slicing, and symbolic execution) to check program
properties given in a form of finite state automata. We have discussed a
synergy of the three methods. Moreover, our experimental results indicate that
the technique can recognize some false positives and some real errors in error
reports produced by other error-detection tools.

%Example of false-negatives/positives -- where it does not work, description of
%where it does not work.

\subsection*{Availability}
\noindent Web pages: \url{http://sf.net/projects/symbiotic/} \\
\noindent Repository (Slicer): \url{http://github.com/jirislaby/LLVMSlicer/} \\
\noindent License: GNU General Public License v2.0

%}}}
